\subsubsection{Subproceso 1: Acreditación y Pre-requisitos}

El subproceso comienza días antes de la fecha de matrícula regular. Por el lado de la universidad, los Departamentos Académicos y las Facultades cierran las actas de notas del semestre anterior sin dejar reclamos pendientes, mientras que las Escuelas Profesionales envían a la DIRCE las actualizaciones por cambios de plan de estudios, convalidaciones y la lista de alumnos egresantes con turno preferencial.

En paralelo, el estudiante genera su orden de pago y cancela el autoseguro estudiantil obligatorio, además de regularizar cualquier deuda de libros o materiales en su Facultad. Si el alumno tiene cursos en riesgo académico, asiste previamente a su sesión de tutoría para que el tutor autorice la activación de su código en el sistema.

Con las notas cerradas y los padrones actualizados, el sistema SIGA-ORCE calcula el promedio ponderado de los dos últimos ciclos regulares de cada estudiante, define su tope máximo de créditos y lo ubica en el grupo y horario de matrícula que le corresponde por orden de mérito. Al mismo tiempo, valida que el pago del autoseguro esté registrado y que no existan bloqueos ni adeudos reportados.

Finalmente, la Intranet DIRCE publica el cronograma de grupos de matrícula de la Facultad y muestra al alumno su ficha de matrícula personalizada junto con el estado de sus requisitos en verde. Con el estudiante habilitado y su turno asignado, concluye el Subproceso 1 y queda listo el ingreso al Subproceso 2.
